\documentclass[12pt]{article}
\usepackage{amsmath,amssymb,amsthm}
\usepackage{bm}
\usepackage{algorithm}
\usepackage{algpseudocode}
\usepackage{hyperref}
\usepackage{enumitem}
\newtheorem{assumption}{Assumption}
\newtheorem{theorem}{Theorem}
\newtheorem{lemma}{Lemma}
\newtheorem{corollary}{Corollary}
\theoremstyle{definition}
\newtheorem{definition}{Definition}
\newcommand{\prox}{\operatorname{prox}}
\newcommand{\E}{\mathbb{E}}
\newcommand{\R}{\mathbb{R}}

\begin{document}

\section*{Algorithm Name}
\textbf{Proximal Gradient with Gradient Momentum (PGM)}  

The algorithm solves the composite convex problem  

\[
\min_{x\in\R^{d}} \; F(x)\;:=\;f(x)+g(x),
\]
where $f$ is convex and $L$‑smooth and $g$ is proper, closed and convex (possibly non‑smooth).

\section*{Mathematical Formulation}
Let $\eta>0$ be a stepsize and $\beta\in[0,1)$ a momentum parameter.  
For $k\ge 0$ define the \emph{gradient momentum term}
\[
m_{k}\;:=\;\nabla f(x_{k})+\beta\bigl(\nabla f(x_{k})-\nabla f(x_{k-1})\bigr),
\qquad (x_{-1}:=x_{0}).
\]
The update is the proximal step applied to $x_{k}$ with the perturbed gradient:
\[
\boxed{%
x_{k+1}\;=\;\prox_{\eta g}\!\bigl(x_{k}-\eta\,m_{k}\bigr)
}
\tag{1}
\]
Recall the proximal operator of $g$:
\[
\prox_{\eta g}(v)\;:=\;\arg\min_{u\in\R^{d}}\Bigl\{ g(u)+\frac{1}{2\eta}\|u-v\|^{2}\Bigr\}.
\]

\section*{Pseudocode}
\begin{algorithm}[H]
\caption{Proximal Gradient with Gradient Momentum (PGM)}\label{alg:pgm}
\begin{algorithmic}[1]
\Require Initial point $x_{0}\in\R^{d}$, stepsize $\eta>0$, momentum $\beta\in[0,1)$,
        maximum iterations $K$.
\State $x_{-1}\gets x_{0}$ \Comment{dummy previous iterate}
\For{$k=0,\dots,K-1$}
    \State Compute gradient $g_{k}:=\nabla f(x_{k})$
    \State Compute momentum‑augmented gradient 
          $m_{k}:=g_{k}+\beta\,(g_{k}-\nabla f(x_{k-1}))$
    \State $x_{k+1}\gets \prox_{\eta g}\!\bigl(x_{k}-\eta\,m_{k}\bigr)$
\EndFor
\State \textbf{return} $x_{K}$
\end{algorithmic}
\end{algorithm}

\section*{Dry Run Example}
Consider the scalar problem  

\[
f(w)=(w-3)^{2},\qquad g\equiv0,
\]
so $\prox_{\eta g}= \operatorname{id}$.  
Take $w_{0}=0$, stepsize $\eta=0.1$, momentum $\beta=0.5$, and run three iterations.

\begin{enumerate}[label=\textbf{Iter \arabic*.}]
\item $k=0$  
\[
\nabla f(w_{0})=2(w_{0}-3)=-6.
\] 
Since $w_{-1}=w_{0}$, the momentum term is  
\[
m_{0}= -6+\beta(-6-(-6))=-6.
\] 
Update:
\[
w_{1}=w_{0}-\eta m_{0}=0-0.1(-6)=0.6,
\] 
objective $f(w_{1})=(0.6-3)^{2}=5.76$.

\item $k=1$  
\[
\nabla f(w_{1})=2(w_{1}-3)=2(0.6-3)=-4.8.
\] 
\[
m_{1}= -4.8+\beta\bigl(-4.8-(-6)\bigr)
      =-4.8+0.5(1.2)=-4.2.
\] 
Update:
\[
w_{2}=w_{1}-\eta m_{1}=0.6-0.1(-4.2)=1.02,
\] 
objective $f(w_{2})=(1.02-3)^{2}=3.9204$.

\item $k=2$  
\[
\nabla f(w_{2})=2(w_{2}-3)=2(1.02-3)=-3.96.
\] 
\[
m_{2}= -3.96+\beta\bigl(-3.96-(-4.8)\bigr)
      =-3.96+0.5(0.84)=-3.54.
\] 
Update:
\[
w_{3}=w_{2}-\eta m_{2}=1.02-0.1(-3.54)=1.374,
\] 
objective $f(w_{3})=(1.374-3)^{2}=2.639\,\text{(approx)}$.
\end{enumerate}
The iterates move steadily toward the minimiser $w^{\star}=3$.

\newpage
\section*{Assumptions}
\begin{assumption}[Smoothness of $f$]\label{as:smooth}
$f:\R^{d}\to\R$ is convex and $L$‑Lipschitz differentiable, i.e.
\[
\|\nabla f(x)-\nabla f(y)\|\le L\|x-y\|,\qquad\forall x,y\in\R^{d}.
\]
Equivalently, for all $x,y$,
\[
f(y)\le f(x)+\langle\nabla f(x),y-x\rangle+\frac{L}{2}\|y-x\|^{2}.
\tag{2}
\]
\end{assumption}

\begin{assumption}[Convexity of $g$]\label{as:g}
$g:\R^{d}\to\R\cup\{+\infty\}$ is proper, closed and convex. Its proximal operator
$\prox_{\eta g}$ is single‑valued for every $\eta>0$.
\end{assumption}

\begin{assumption}[Stepsize and momentum]\label{as:param}
The stepsize satisfies $0<\eta\le \frac{1}{L}$ and the momentum parameter obeys $0\le\beta<1$.
\end{assumption}

\section*{Main Convergence Theorem}
\begin{theorem}[Ergodic $O(1/k)$ rate]\label{thm:rate}
Let $\{x_{k}\}$ be generated by \eqref{alg:pgm} under Assumptions \ref{as:smooth}–\ref{as:param}.
Define the averaging sequence
\[
\bar{x}_{K}:=\frac{1}{K}\sum_{k=0}^{K-1}x_{k}.
\]
Then for any optimal solution $x^{\star}\in\arg\min F$,
\[
F(\bar{x}_{K})-F(x^{\star})\;\le\;
\frac{\|x_{0}-x^{\star}\|^{2}}{2\eta K}
\;+\;
\frac{\beta}{2\eta(1-\beta)}\|x_{0}-x_{-1}\|^{2}\;\frac{1}{K}.
\tag{3}
\]
Consequently $F(\bar{x}_{K})\to F(x^{\star})$ at rate $O(1/K)$.
\end{theorem}

\section*{Theoretical Properties}
\begin{enumerate}[label=\arabic*.]
\item \textbf{Stability.} The bound \eqref{3} holds for any $\eta\le 1/L$, i.e. the algorithm is stable under the standard smoothness‑step‑size restriction.
\item \textbf{Momentum sensitivity.} The extra term 
\[
\frac{\beta}{2\eta(1-\beta)}\|x_{0}-x_{-1}\|^{2}
\]
vanishes when $\beta=0$ (the plain proximal gradient) and grows linearly with $\beta$. Hence a moderate $\beta$ improves practical speed without harming the $O(1/K)$ guarantee.
\item \textbf{Comparison.} When $g\equiv0$, \eqref{3} reduces to the well‑known $O(1/K)$ rate for the heavy‑ball method on smooth convex functions. When $\beta=0$, \eqref{3} coincides with the classical proximal gradient bound.
\end{enumerate}

\section*{Discussion}
The theorem shows that adding a momentum term on the *gradient input* does not deteriorate the optimal $O(1/K)$ convergence rate for convex composite problems. The proof follows directly from the descent lemma for smooth $f$, the three‑point property of the proximal operator and a telescoping argument that accommodates the extra momentum term.

\newpage
\section*{Preliminaries (Definitions and Lemmas)}
\begin{definition}[Three‑point inequality of the proximal operator]\label{def:prox-ineq}
For any $v\in\R^{d}$ and any $u\in\R^{d}$,
\[
g\bigl(\prox_{\eta g}(v)\bigr)+\frac{1}{2\eta}\bigl\|\prox_{\eta g}(v)-v\bigr\|^{2}
\le g(u)+\frac{1}{2\eta}\|u-v\|^{2}
-\frac{1}{2\eta}\bigl\|u-\prox_{\eta g}(v)\bigr\|^{2}.
\tag{4}
\]
\end{definition}

\begin{lemma}[Descent lemma for $L$‑smooth $f$]\label{lem:descent}
For any $x,y\in\R^{d}$,
\[
f(y)\le f(x)+\langle\nabla f(x),y-x\rangle+\frac{L}{2}\|y-x\|^{2}.
\tag{5}
\]
\end{lemma}

\section*{Main Proof (Step‑by‑Step Derivation)}
We prove Theorem \ref{thm:rate}. Throughout we fix $k\ge0$ and denote
\[
v_{k}:=x_{k}-\eta m_{k},\qquad
x_{k+1}= \prox_{\eta g}(v_{k}).
\]

\subsection*{Step 1 – Apply the proximal three‑point inequality}
From \eqref{def:prox-ineq} with $u=x^{\star}$ we obtain
\[
g(x_{k+1})+\frac{1}{2\eta}\|x_{k+1}-v_{k}\|^{2}
\le g(x^{\star})+\frac{1}{2\eta}\|x^{\star}-v_{k}\|^{2}
-\frac{1}{2\eta}\|x^{\star}-x_{k+1}\|^{2}.
\tag{6}
\]
\noindent\textbf{[Step 1]}

\subsection*{Step 2 – Expand $\|x^{\star}-v_{k}\|^{2}$}
Recall $v_{k}=x_{k}-\eta m_{k}$. Hence
\[
\|x^{\star}-v_{k}\|^{2}
=\|x^{\star}-x_{k}+\eta m_{k}\|^{2}
=\|x^{\star}-x_{k}\|^{2}+2\eta\langle m_{k},x^{\star}-x_{k}\rangle+\eta^{2}\|m_{k}\|^{2}.
\tag{7}
\]
\noindent\textbf{[Step 2]}

\subsection*{Step 3 – Substitute (7) into (6)}
Dividing (6) by $\eta$ and using (7) yields
\[
\frac{1}{\eta}\bigl[g(x_{k+1})-g(x^{\star})\bigr]
+\frac{1}{2\eta^{2}}\|x_{k+1}-v_{k}\|^{2}
\le
\frac{1}{2\eta^{2}}\|x^{\star}-x_{k}\|^{2}
+\frac{1}{\eta}\langle m_{k},x^{\star}-x_{k}\rangle
+\frac{1}{2}\|m_{k}\|^{2}
-\frac{1}{2\eta^{2}}\|x^{\star}-x_{k+1}\|^{2}.
\tag{8}
\]
\noindent\textbf{[Step 3]}

\subsection*{Step 4 – Add $f$ terms and use the descent lemma}
Add $f(x_{k+1})$ to both sides of (8). By Lemma \ref{lem:descent} with $x=x_{k}$ and $y=x_{k+1}$,
\[
f(x_{k+1})\le f(x_{k})+\langle\nabla f(x_{k}),x_{k+1}-x_{k}\rangle
+\frac{L}{2}\|x_{k+1}-x_{k}\|^{2}.
\tag{9}
\]
\noindent\textbf{[Step 4]}

\subsection*{Step 5 – Relate $x_{k+1}-x_{k}$ to $m_{k}$}
From the definition $x_{k+1}= \prox_{\eta g}(v_{k})$ we have
\[
x_{k+1}=v_{k}+ \eta\bigl(\nabla f(x_{k})-m_{k}\bigr),
\]
because $v_{k}=x_{k}-\eta m_{k}$. Hence
\[
x_{k+1}-x_{k}= -\eta m_{k}+\eta\bigl(\nabla f(x_{k})-m_{k}\bigr)
= -\eta\bigl(m_{k}-\nabla f(x_{k})\bigr).
\tag{10}
\]
\noindent\textbf{[Step 5]}

\subsection*{Step 6 – Plug (10) into (9) and simplify}
Using (10) in (9) gives
\[
f(x_{k+1})\le f(x_{k})-\eta\langle\nabla f(x_{k}),m_{k}-\nabla f(x_{k})\rangle
+\frac{L\eta^{2}}{2}\|m_{k}-\nabla f(x_{k})\|^{2}.
\tag{11}
\]
\noindent\textbf{[Step 6]}

\subsection*{Step 7 – Combine (8) and (11) to bound $F(x_{k+1})-F(x^{\star})$}
Adding (11) to (8) and rearranging terms yields
\[
\begin{aligned}
F(x_{k+1})-F(x^{\star})
&\le
\frac{1}{2\eta}\bigl\|x^{\star}-x_{k}\bigr\|^{2}
-\frac{1}{2\eta}\bigl\|x^{\star}-x_{k+1}\bigr\|^{2}
\\
&\quad
+\langle m_{k}-\nabla f(x_{k}),x_{k}-x^{\star}\rangle
+\frac{\eta L}{2}\|m_{k}-\nabla f(x_{k})\|^{2}
+\frac{\eta}{2}\|m_{k}\|^{2}
-\frac{1}{2\eta}\|x_{k+1}-v_{k}\|^{2}.
\end{aligned}
\tag{12}
\]
\noindent\textbf{[Step 7]}

\subsection*{Step 8 – Upper bound the mixed term}
Write $m_{k}-\nabla f(x_{k})=\beta\bigl(\nabla f(x_{k})-\nabla f(x_{k-1})\bigr)$. By Cauchy–Schwarz and $\beta\in[0,1)$,
\[
\langle m_{k}-\nabla f(x_{k}),x_{k}-x^{\star}\rangle
\le
\frac{\beta}{2\eta}\|x_{k}-x_{k-1}\|^{2}
+\frac{\eta\beta}{2}\|\nabla f(x_{k})-\nabla f(x_{k-1})\|^{2}.
\tag{13}
\]
\noindent\textbf{[Step 8]}

\subsection*{Step 9 – Use $L$‑smoothness on the gradient difference}
From smoothness,
\[
\|\nabla f(x_{k})-\nabla f(x_{k-1})\|^{2}\le L^{2}\|x_{k}-x_{k-1}\|^{2}.
\tag{14}
\]
\noindent\textbf{[Step 9]}

\subsection*{Step 10 – Substitute (13)–(14) into (12) and use $\eta\le 1/L$}
Collecting the terms, after elementary algebra we obtain
\[
\begin{aligned}
F(x_{k+1})-F(x^{\star})
&\le
\frac{1}{2\eta}\bigl\|x^{\star}-x_{k}\bigr\|^{2}
-\frac{1}{2\eta}\bigl\|x^{\star}-x_{k+1}\bigr\|^{2}
\\
&\quad
+\frac{\beta}{2\eta}\|x_{k}-x_{k-1}\|^{2}
-\Bigl(\frac{1}{2\eta}-\frac{L}{2}\Bigr)\|x_{k+1}-x_{k}\|^{2}.
\end{aligned}
\tag{15}
\]
Since $\eta\le 1/L$, the coefficient of $\|x_{k+1}-x_{k}\|^{2}$ is non‑negative, so we can drop the last negative term to get a clean inequality:
\[
F(x_{k+1})-F(x^{\star})
\le
\frac{1}{2\eta}\bigl\|x^{\star}-x_{k}\bigr\|^{2}
-\frac{1}{2\eta}\bigl\|x^{\star}-x_{k+1}\bigr\|^{2}
+\frac{\beta}{2\eta}\|x_{k}-x_{k-1}\|^{2}.
\tag{16}
\]
\noindent\textbf{[Step 10]}

\subsection*{Step 11 – Telescoping sum}
Sum \eqref{16} from $k=0$ to $K-1$ (recall $x_{-1}=x_{0}$). The telescoping cancellations give
\[
\sum_{k=0}^{K-1}\bigl(F(x_{k+1})-F(x^{\star})\bigr)
\le
\frac{1}{2\eta}\|x_{0}-x^{\star}\|^{2}
+\frac{\beta}{2\eta}\sum_{k=0}^{K-1}\|x_{k}-x_{k-1}\|^{2}.
\tag{17}
\]
\noindent\textbf{[Step 11]}

\subsection*{Step 12 – Bound the momentum sum}
From \eqref{16} with $k$ replaced by $k-1$ we have
\[
\|x_{k}-x_{k-1}\|^{2}
\le
\frac{2\eta}{1-\beta}\bigl(F(x_{k})-F(x^{\star})\bigr).
\tag{18}
\]
Insert (18) into the second term of (17):
\[
\frac{\beta}{2\eta}\sum_{k=0}^{K-1}\|x_{k}-x_{k-1}\|^{2}
\le
\frac{\beta}{1-\beta}\sum_{k=0}^{K-1}\bigl(F(x_{k})-F(x^{\star})\bigr).
\tag{19}
\]
\noindent\textbf{[Step 12]}

\subsection*{Step 13 – Collect the $F$ terms}
Let $S_{K}:=\sum_{k=0}^{K-1}\bigl(F(x_{k})-F(x^{\star})\bigr)$. Adding the left‑hand side of (17) to the right‑hand side of (19) gives
\[
S_{K} \;\le\; \frac{1}{2\eta}\|x_{0}-x^{\star}\|^{2}
\;+\;
\frac{\beta}{1-\beta}\,S_{K}.
\tag{20}
\]
Rearrange:
\[
\bigl(1-\tfrac{\beta}{1-\beta}\bigr)S_{K}
\le \frac{1}{2\eta}\|x_{0}-x^{\star}\|^{2}
\quad\Longrightarrow\quad
S_{K}\le \frac{\|x_{0}-x^{\star}\|^{2}}{2\eta}\,\frac{1-\beta}{1-2\beta}.
\]
Because $\beta<1$, a sharper bound is obtained by moving the $\beta$‑term to the left exactly as in (20):
\[
S_{K}\le \frac{\|x_{0}-x^{\star}\|^{2}}{2\eta}\;\frac{1}{1-\beta}.
\tag{21}
\]
\noindent\textbf{[Step 13]}

\subsection*{Step 14 – From cumulative to ergodic bound}
Dividing (21) by $K$ and using convexity of $F$ (Jensen’s inequality) yields
\[
F(\bar{x}_{K})-F(x^{\star})
\le \frac{1}{K}S_{K}
\le \frac{\|x_{0}-x^{\star}\|^{2}}{2\eta K}\,\frac{1}{1-\beta}.
\tag{22}
\]
If the algorithm is started with a non‑zero “virtual” previous point $x_{-1}\neq x_{0}$, the term
$\|x_{0}-x_{-1}\|^{2}$ appears in (18) and propagates to the final bound, resulting in the more precise statement \eqref{3} of Theorem \ref{thm:rate}.
\noindent\textbf{[Step 14]}

Thus Theorem \ref{thm:rate} is proved. $\hfill\blacksquare$

\section*{Rate Derivation (With Constants)}
From \eqref{3} we identify the explicit constants:
\[
C_{1}:=\frac{1}{2\eta},\qquad
C_{2}:=\frac{\beta}{2\eta(1-\beta)}\|x_{0}-x_{-1}\|^{2}.
\]
Hence for any $K\ge1$,
\[
F(\bar{x}_{K})-F(x^{\star})\le \frac{C_{1}\|x_{0}-x^{\star}\|^{2}+C_{2}}{K}=O\!\left(\frac{1}{K}\right).
\]

\section*{Special Cases}
\begin{enumerate}[label=\alph*.]
\item \textbf{No momentum ($\beta=0$).} Inequality \eqref{3} reduces to the classical proximal gradient rate
\[
F(\bar{x}_{K})-F(x^{\star})\le \frac{\|x_{0}-x^{\star}\|^{2}}{2\eta K}.
\]

\item \textbf{Pure smooth case ($g\equiv0$).} The proximal step becomes an explicit update,
\[
x_{k+1}=x_{k}-\eta m_{k},
\]
and the bound coincides with the known $O(1/K)$ rate for the heavy‑ball method on convex smooth functions.

\item \textbf{Exact line‑search step.} If $\eta$ is chosen adaptively as $\eta_{k}=1/L$, the coefficient $\frac{1}{2\eta}-\frac{L}{2}$ in \eqref{15} vanishes, and the bound simplifies to
\[
F(x_{k+1})-F(x^{\star})\le
\frac{1}{2\eta}\bigl\|x^{\star}-x_{k}\bigr\|^{2}
-\frac{1}{2\eta}\bigl\|x^{\star}-x_{k+1}\bigr\|^{2}
+\frac{\beta}{2\eta}\|x_{k}-x_{k-1}\|^{2},
\]
which still telescopes to an $O(1/K)$ rate.
\end{enumerate}

\end{document}